
La hibridación de ADN es el proceso mediante el cual dos cadenas de ADN o ARN se unen con secuencias complementarias mediante apareamiento de bases.

Algunos tipos son:

\begin{itemize}
    \item Northern blot, para ARN
    \item Southern blot, para ADN
    \item In Situ
    \item Hibridación fluorescente in situ (Fluorescent In Situ Hibridation o FISH)
    \item Microarreglos
\end{itemize}

\section{Desnaturalización}

Primero, se inicia separando las cadenas de ADN/ARN, esto mediante varios mecanismos, ya sea el aumento de temperatura, la inducción de cambios  en el pH o el uso de agentes químicos.

\section{Hibridación}

Luego se introducen las sondas complementarias. Estas sondas son fragmentos de ADN o ARN los cuales son diseñados para reconocer mediante complementariedad una secuencia específica de interés.

\section{Detección}

La detección se logra al marcar las sondas complementarias para luego medir tal marcaje. El marcaje de estas sondas puede realizarse con varios métodos, ya sean radioactivos, como el uso de Plomo 32 (\ce{^{32}P}), o no radioactivos, como la biotina la dioxigenina, el uso de fluoróforos o también enzimas-

\section{Stringencia}

Este parámetro mide el grado de rigurosidad de las condicines utilizadas para permitir o impedir la formación de híbridos imperfectamente complementarios. Un valor alto de strigencia indica mayor especificidad, mientras que unvalor bajo de strigencia refleja una menor especificidad

Este valor puede ser afectado por el manejo de la temperatura durante el proceso, la concentración de sales en la muestra, la longitud de la sonda complementaria, así como el contenido de Citosina/Guanina (CG) en el ADN de interés\footnote{La unión CG es más fuerte (3 puentes de hidrógeno) que la AT (2 puentes de hidrógeno), por lo que un mayor contenido de CG hace el paso la desnaturalización más dificil.}